\documentclass{article}

% paquete para imágenes y figuras del proyecto
\usepackage{graphicx}
% paquete para la bibliografía
\usepackage[backend=bibtex,sorting=none]{biblatex}
\bibliography{MNIG_to_FMMSP_Tatiana_Porras}
% paquete de símbolos matemáticos
\usepackage{amsfonts}

% metadata del proyecto
\title{MNIG to FMMSP Tatiana Porras}
\date{2020-03-13}
\author{Tatiana Porras}

% definición de variables
\def\direcM{LINDSAY ÁLVAREZ POMAR}
\def\direc{Lindsay Álvarez Pomar}
\def\algmod{MNIG\_to\_FMMSP} % algmod es por algoritmo modelo
\makeatletter
\let\dateVar\@date
\makeatother

% tesis:
\begin{document}

% página de título
\begin{titlepage}
    \begin{center}
    ALGORITMO VORAZ ITERATIVO CON MULTI-VECINDAD APLICADO AL PROBLEMA DE SECUENCIACIÓN DIFUSO MULTIPRODUCTO Y MULTIETAPAS
    \par \bigskip
    Anteproyecto de Grado
    \par \vspace{5cm}
    Tatiana Porras Cortes
    \par \medskip
    Correo \hspace{1cm} tatiporras96@gmail.com
    \par \bigskip
    LINDSAY ÁLVAREZ POMAR \par Director del trabajo de grado
    \par \vspace{1cm}
    \begin{figure}[h!]
        \begin{center}
        \includegraphics[width=6cm]{EscudoUD1.png}
        \end{center}
    \end{figure}
    \par UNIVERSIDAD DISTRITAL FRANCISCO JOSÉ DE CALDAS
    \par FACULTAD DE INGENIERÍA
    \par PROYECTO CURRICULAR DE INGENIERÍA INDUSTRIAL
    \par Bogotá D.C., Colombia. \dateVar
    \end{center}
\end{titlepage}

\renewcommand{\contentsname}{Tabla de Contenido}

\tableofcontents

\setcounter{section}{-1}

% secciones del proyecto
\section{TÍTULO}

"ALGORITMO VORAZ ITERATIVO CON MULTI-VECINDAD APLICADO AL PROBLEMA DE
SECUENCIACIÓN DIFUSO MULTIPRODUCTO Y \linebreak MULTIETAPAS"

\newpage

\section{RESUMEN}

El Algoritmo Voraz Iterativo con Multi-Vecindad aplicado al Problema de
Secuenciación Difuso Multiproducto y Multietapas (\algmod\ por sus siglas
en inglés) no ha sido tratado en la literatura científica internacional.
En este trabajo se explora esta combinación nueva. Para ello, se modificará
el algoritmo MNIG de modo que pueda ser utilizado para resolver el modelo FMMSP.

\vspace{\baselineskip}
Al final se reportarán los resultados del algoritmo aplicado a varias instancias 
del problema, y se compara dichos resultados con los de otros cuatro algoritmos 
que han sido aplicados al modelo FMMSP en la literatura.

\vspace{\baselineskip}
\textbf{Palabras Clave}: Algoritmo Voraz Iterativo, Multi-vecindad,
Difuso, Secuenciación, Multiproducto, Multietapas

\subsection{ABSTRACT}

The Multi-neighborhood Iterated Greedy algorithm applied to the Fuzzy 
Multiproduct Multistage Scheduling Problem (from now on \algmod)
has not been treated in the international scientific literature. This
work explores this new combination. For that, the MNIG algorithm will be
modified so that it can be used to solve the FMMSP model.

\vspace{\baselineskip}
At the end, the results of the algorithm will be reported, after applying it to 
several instances of the problem, and the results will be compared to those of 
another four algorithms that have been applied to the FMMSP model in the 
literature. 

\vspace{\baselineskip}
\par \textbf{Keywords}: Iterated Greedy, Multi-neighborhood, Scheduling, Fuzzy,
Multiproduct, Multistage

\section{INTRODUCCIÓN}

Mediante el presente trabajo de tesis de pregrado se pretende hacer un aporte
aunque pequeño al conocimiento. En la ingeniería industrial existen diversas
áreas en las que se puede hacer un aporte de este tipo. Existe el área de
investigación de operaciones, el área de gestión, mercadeo, higiene industrial,
seguridad y salud en el trabajo, ergonomía, etcétera. 

\vspace{\baselineskip}
Un aporte de nuevo conocimiento en pregrado es un aporte pequeño y muy
específico, detallado. Por ello se selecciona una de las áreas de conocimiento
de la ingeniería industrial y dentro de esa área se elige un tema en particular.
De ese tema elegido se trabaja un detalle que no haya sido estudiado con
anterioridad. 

\vspace{\baselineskip}
En el presente trabajo se eligió el tema de secuenciación, mejor conocido por
su nombre en inglés como "scheduling". Dentro de este tema se revisó el estado
del arte, y se encontró un modelo de scheduling que ha sido poco estudiado, y
por aparte se encontró un algoritmo que nunca ha sido aplicado al modelo, pero
que podría ser aplicado. El modelo encontrado \autocite{modFMMSP} es llamado 
FMMSP que significa Fuzzy Multiproduct Multistage Scheduling Problem, que en 
español se podría traducir como Problema de Secuenciación Difuso Multiproducto 
y Multietapas. Es un modelo interesante de ser estudiado pues incluye números 
difusos, múltiples productos, y múltiples etapas. El algoritmo encontrado 
\autocite{algMNIG} es llamado MNIG que significa Multi-Neighborhood Iterated
Greedy algorithm. En español se traduce como algoritmo Voraz Iterativo con 
Multi-Vecindad. Se ha encontrado que este algoritmo sirve para resolver 
problemas de scheduling \autocite{algMNIG}. 

\vspace{\baselineskip}
Lo novedoso resulta en que dicho algoritmo jamás ha sido aplicado al modelo, a 
tal punto que para realizar este trabajo se hará necesario modificar el 
algoritmo, pues el algoritmo original no se puede aplicar directamente al 
modelo. No existen pruebas de que este algoritmo sea apropiado para este modelo.
Con el presente trabajo se propone comprobar que tan bueno es el algoritmo MNIG
aplicado al modelo FMMSP, al compararlo con otros cuatro algoritmos que ya han
sido aplicados al modelo FMMSP \autocite{modFMMSP}.

\section{PLANTEAMIENTO DEL PROBLEMA}

El problema que ataca esta tesis es explorar nuevas posibilidades de encontrar
secuencias en las que se puedan realizar los trabajos en un sistema de
producción del tipo FMMSP, de modo tal que dichas secuencias de trabajos
tengan un tiempo de terminación máximo lo más bajo posible.

\vspace{\baselineskip}
Siendo que el espacio de búsqueda del modelo FMMSP para problemas de tamaño
realista es muy grande como para buscar de manera exhaustiva la secuencia con
el menor de los tiempos de terminación máximo (es decir probando cada secuencia
posible por separado), se diseñan algoritmos que permitan buscar en regiones
promisorias del espacio de búsqueda. De esta manera, se hace innecesario 
revisar todo el espacio de búsqueda, y los algoritmos dan buenos resultados,
o sea, con bajos tiempos de terminación máximo.

\vspace{\baselineskip}
El algoritmo MNIG aplicado al modelo FMMSP tiene la posibilidad de competir
en el encuentro de buenas soluciones. Aplicando el algoritmo MNIG un sistema
productivo tipo FMMSP podría encontrar secuencias de trabajos con bajos tiempos
de terminación máximo. Sin embargo esta posibilidad aún no está comprobada, de
eso se trata esta tesis.

\subsection{Elementos del Problema}

Según lo anterior, los elementos del problema son:

\vspace{\baselineskip}
\begin{itemize}
    \item[\maltese] Se requieren secuencias con bajos tiempos de terminación en
    el modelo FMMSP.
    \item[\maltese] El algoritmo MNIG tiene potencial para encontrar dichas
    secuencias en un tiempo prudente. Los resultados se pueden comparar a otros
    cuatro algoritmos que ya han sido usados en el modelo FMMSP
\end{itemize}

\subsection{Pregunta del Problema}

¿Puede aplicarse el algoritmo MNIG al modelo FMMSP para obtener resultados
similares a los obtenidos mediante los otros cuatro algoritmos ya probados?

\subsection{Subpreguntas del Problema}

\indent\indent
¿Puede implementarse el algoritmo MNIG aplicado al modelo FMMSP en un lenguaje
de programación como Python?

¿Puede correr el algoritmo MNIG aplicado al modelo FMMSP e imprimir o devolver
una secuencia de trabajos factible en ese modelo?

¿Es comparablemente bueno el algoritmo MNIG aplicado al modelo FMMSP respecto
a los otros cuatro algoritmos que ya han sido aplicados a ese modelo?

\section{OBJETIVOS}

\subsection{Objetivo General}

\begin{itemize}
    \item Aplicar el algoritmo MNIG al modelo FMMSP
\end{itemize}

\subsection{Objetivos Específicos}

\begin{itemize}
    \item Implementar el algoritmo MNIG aplicado al modelo FMMSP en el
    \linebreak lenguaje de programación Python
    \item Correr el algoritmo MNIG en instancias del modelo FMMSP y
    obtener secuencias factibles
    \item Comparar el algoritmo MNIG aplicado al modelo FMMSP con los otros
    cuatro algoritmos que se han aplicado al modelo FMMSP. La comparación se
    hará con el tiempo de terminación máximo de la secuencia resultante
\end{itemize}

\section{JUSTIFICACIÓN}



\section{ALCANCES}



\section{MARCO DE REFERENCIA}



\section{HIPÓTESIS}



\section{DISEÑO METODOLÓGICO}



\section{CRONOGRAMA}



\section{PRODUCTOS DEL PROYECTO}



\section{BIBLIOGRAFÍA}

\printbibliography[heading=none]

\end{document}